\section*{Your turn}

Check out tag \code{ch13} of the companion repository and run \code{npm ci}
once. Exercises 1, 2, 5 and 10 need the database, the seven subgraphs and the
router. Exercises 3 and 4 need all that and a second checkout at tag \code{ch12},
because what they measure is the defect this chapter repaired. Exercises 6, 7 and
9 need node and nothing started at all. Exercise 8 needs two more processes and a
router of its own.

There are seven services now, plus the router:

\begin{minted}[fontsize=\footnotesize,breaklines]{text}
docker compose up -d --build mosaic-db mosaic-catalog mosaic-pricing \
    mosaic-inventory mosaic-accounts mosaic-reviews mosaic-ordering \
    mosaic-nodes mosaic-router
docker compose ps
\end{minted}

Nine containers, all healthy. \code{mosaic-nodes} is the one with no database
beside it in the table; if it is the only one that will not start, something else
has 5107.

\begin{enumerate}
  \item \textbf{Follow one identifier.} Take a product identifier out of
  \code{browseProducts}, and send this to port 3002 with the two query-plan
  headers chapter~\ref{ch:enter-the-router} used:

  \begin{minted}[fontsize=\footnotesize,breaklines]{text}
X-WG-Include-Query-Plan: true
X-WG-Skip-Loader: true
  \end{minted}

  \begin{minted}[fontsize=\footnotesize,breaklines]{graphql}
{ node(id: "...") { ... on Product {
    title price { amount } availableQuantity } } }
  \end{minted}

  Count the fetches in \code{extensions.queryPlan}, then look at what the first
  one asked the node service for. It is two fields and nothing else, and the
  whole rest of the answer is keyed on them.

  \item \textbf{Ask for something that is not there.} Base64-encode
  \code{Product:} followed by sixteen zero bytes and send that to \code{node},
  selecting \code{title}. Note which service the error comes from, and which
  service it is actually about. Then send the same identifier selecting only
  \code{\_\_typename} and \code{id}: it answers \code{Product} and hands the
  identifier straight back, because nothing on that path ever had to find a row.
  Decide whether you would rather \code{node} answered null instead, and what the
  node service would have to do to manage it.

  \item \textbf{Break the paging on purpose.} Ask the router for
  \code{browseProducts(first: 2)} selecting \code{cursor} and
  \code{node \{ title \}}, and base64-decode the second cursor. Now do the same
  thing at tag \code{ch12}, in a second checkout, and decode it again. One of
  them ends in an all-zero \code{Guid}. Use each cursor as \code{after:} and see
  which page repeats a row.

  \item \textbf{Then hide the bug the way this book did for eight chapters.}
  Still at \code{ch12}, add \code{price \{ amount \}} to the same query and
  decode the cursor again. It is correct. Write down, before reading
  section~\ref{sec:ch13-the-page-that-fits} again, why one more field from
  another service repaired a cursor Catalog built.

  \item \textbf{Ask for the page that cannot exist.} Send
  \code{browseProducts(first: 3, order: [\{ price: \{ amount: ASC \} \}])}.
  Read the error and note where it came from: no subgraph was asked anything.
  Then look at \code{ProductSortInput} through the router's introspection and
  satisfy yourself that no directive would have changed the list.

  \item \textbf{Make the composer merge one enum two ways.} Print these in turn:

  \begin{minted}[fontsize=\footnotesize,breaklines]{text}
node scripts/modeling-cases.mjs --print enum-in-output-only-is-unioned
node scripts/modeling-cases.mjs --print enum-in-input-only-is-intersected
  \end{minted}

  The two cases declare the same two sets of members and get opposite answers.
  Read both case definitions in \code{scripts/modeling-cases.mjs} and find the
  two differences between them, then work out which of the two is doing the
  work. Then decide which merge you would rather discover in production, and
  what would have told you.

  \item \textbf{Watch a composer stop being a composer.} Print the case named
  for an implementation without a key:

  \begin{minted}[fontsize=\footnotesize,breaklines]{text}
node scripts/modeling-cases.mjs --print interface-object-implementation-without-a-key
  \end{minted}

  Then read the case in \code{scripts/modeling-cases.mjs} and find the two edits
  it makes to \code{schema/samples/interface-object-library.graphql}. One of them
  is the mistake and the other is the bookkeeping that follows from it. Consider
  what a build pipeline does with that exit code, and what a reviewer would have
  to know to catch the one-directive cause.

  \item \textbf{Give a type a field nobody declared on it.} Start the two
  interface-object subgraphs, compose them, and bring up the router on 3004:

  \begin{minted}[fontsize=\footnotesize,breaklines]{text}
dotnet run --project samples/interface-object/Mosaic.Sample.InterfaceObject.Library
dotnet run --project samples/interface-object/Mosaic.Sample.InterfaceObject.Ratings
npx wgc router compose -i samples/interface-object/graph.yaml \
                       -o samples/interface-object/supergraph.json
docker compose --profile interfaces up -d interfaces-router
  \end{minted}

  Ask 3004 for \code{\{ library \{ \_\_typename title averageRating \} \}}, then
  ask 5208 for \code{\{ resolutionLog \}}. Compare what the ratings subgraph was
  asked with what the client asked for. Then add a third implementation to
  \code{Library.cs}, a \code{Podcast} with its own key, recompose, and see what
  it costs the ratings subgraph to support it.

  \item \textbf{Find out what a scalar is worth.} In a scratch copy of
  \code{schema/ordering.graphql}, change \code{Decimal}'s \code{@specifiedBy}
  url to anything you like and compose the seven by hand. Then change
  \code{Money.currency} to \code{String} in the same copy and compose again.
  Both succeed. Write down which of the two you would expect a reviewer to catch
  in a pull request, and then go and look at how the two edits differ in a diff.

  \item \textbf{Verify it in Postman.} The collection for this chapter is
  \code{postman/mosaic-nodes}, seven requests against the router:

  \begin{minted}[fontsize=\footnotesize,breaklines]{text}
npx newman run postman/mosaic-nodes.postman_collection.json \
    -e postman/mosaic-nodes.local.postman_environment.json
  \end{minted}

  Twenty-seven assertions, all green at \code{ch13}. Now run the same collection
  against a \code{ch12} checkout with its own graph up, and sort the failures
  into two piles. One pile is requests that fail because a field is not in the
  schema. The other is the last two requests, which fail with the schema
  unchanged, because a page repeats a row. Decide which pile a
  schema-comparison tool would have warned you about, and what you would have
  had to run to catch the other.
\end{enumerate}

Exercise 4 is the one to do properly. Everything else in this chapter is a
property of federation; that one is a property of testing anything behind a
layer that rewrites the request. The query the client wrote and the query the
subgraph answered were different documents, so the subgraph was never tested
against the thing that was broken.
Chapter~\ref{ch:automated-tests} is where a test suite gets told about that
distinction.
